\section{学历的两面：它度量的是什么}

把前两章放在一起看，会得出一个让人舒服、也让人清醒的判断：

\key{学历度量的大多是「你拿到了多少外部脚手架」，而不是「你有多大的内在能力」。}

同样是广度算法的大脑，放进有老师推、有家长托、有资源垫的环境里，能走出高学历；放进没人管、没人信的环境里，可能连高中都读不完。这不是能力之差，是环境之差。

\subsection{「正常人被极大高估了」}

原视频在一期叫「ADHD 算是病吗」的视频里，讲了一个颠覆性的观点（BV1PLXQBUEns）：\term{和 ADHD 对比的时候，正常人往往被极大高估了}。

家长最怕的场景是：「我孩子一个晚上一个单词都搞不定，隔壁小孩一天一百个，这差距不是越来越大？」但原视频说：\term{这个线性积累模型本身就是错的}。知识不是线性累积的——按部就班学了多少年、学得一般的人，多的是。十年后，两个孩子很可能都在中游。

这有两层意思，对学历光谱的两端都成立：

\begin{checklist}
  \item \textbf{对学历中下层}：你当年「追不上别人」的焦虑，可能从头到尾建立在错误的模型上。按部就班不是积累的唯一方式，甚至不是大多数人真正掌握知识的方式；
  \item \textbf{对学历上层}：也别太得意。高学历不能证明你的系统健康，只能证明环境替你做了一部分执行。脚手架一撤，谁裸奔谁知道。
\end{checklist}

\subsection{同一颗大脑，两种剧本}

把两章的画像放在同一张表里，差异一目了然：

\begin{table}[h]
\centering\small
\begin{tabular}{>{\raggedright\arraybackslash}p{2.6cm} p{6.1cm} p{6.1cm}}
\toprule
 & \textbf{学历上层} & \textbf{学历中下层} \\
\midrule
上学体验 & 成绩还行，靠聪明和被动注意力混过 & 反复失败，被归因为「懒 / 笨 / 态度问题」 \\
自我概念 & 「我很强，但不能暴露我很乱」 & 「我不行，做不到是我的宿命」 \\
主要燃料 & 死线肾上腺素 + 超专注红利 & 即时刺激 + 兴趣上头 \\
崩塌方式 & 燃尽循环、越成功越空心 & 回避一切、机会滚雪球式流失 \\
自救重点 & 先建系统，停止靠肾上腺素续命 & 先修自我概念，攒「我能做到」的证据 \\
\bottomrule
\end{tabular}
\end{table}

\subsection{中间态：大多数人都在这条光谱上}

当然，现实里绝大多数人不在两个极端，而在中间：学历不低也不高，症状不重也不轻，时而上头时而下沉。这张光谱的意义不是让你对号入座去站队，而是让你\term{认出自己的燃料和崩塌方式属于哪一型，然后对症下药}。

而且两型的经验是可以互相借用的：

\begin{itemize}
  \item 学历上层可以借中下层的「放下」：允许自己蒙混过关、允许自己放弃不重要的事，先混过去再重新产生兴趣；
  \item 学历中下层可以借上层的「系统」：用恰当场、程式性计划这些方法，给自己搭一个没人替你搭的脚手架。
\end{itemize}

\subsection{三个可干预的变量}

最后，把这张光谱的时间轴拉直，会发现决定一个人滑向哪端的关键变量只有三个，而\term{这三个全部可以后天干预}：

\begin{checklist}
  \item \textbf{自我概念}——可以重建。从「我不行」改写成「我这套系统需要不同的运行方式」；
  \item \textbf{环境脚手架}——可以自建。没人替你做执行，就用恰当场、环境布局、外置大脑替你做；
  \item \textbf{早期经验}——可以重写叙事。过去无法改变，但「我当年为什么失败」的解释可以重写：不是我不行，是我一直缺脚手架。
\end{checklist}

\key{剧本不是定死的。同一个算法，过去把你写成了某一型；从现在起，你可以亲手改写这一章。}
